% ============================================================================
% Lettre de motivation Bell 2027 - version 9 - corrections personnelles
% Style du CV : Carlito, bleu #263068, rouge #E6332A.
% Sans illustrations en marge droite. Compiler avec XeLaTeX.
% ============================================================================
\documentclass[10pt,a4paper]{article}
\usepackage[a4paper,top=1.36cm,bottom=1.30cm,left=1.84cm,right=1.84cm]{geometry}
\usepackage{fontspec}
\setmainfont{Carlito}[BoldFont={Carlito Bold},ItalicFont={Carlito Italic},Ligatures=NoCommon]
\usepackage[french]{babel}
\usepackage{microtype}
\usepackage{xcolor}
\usepackage{setspace}
\usepackage{ragged2e}
\usepackage[hidelinks]{hyperref}
\definecolor{bleu}{HTML}{263068}
\definecolor{rouge}{HTML}{E6332A}
\definecolor{gris}{HTML}{9EA6B6}
\definecolor{anthracite}{HTML}{222936}
\definecolor{secondaire}{HTML}{475269}
\color{anthracite}
\hypersetup{pdftitle={Lettre de motivation - Bell 2027 - Youssef EL KATTOUFI},pdfauthor={Youssef EL KATTOUFI}}
\pagestyle{empty}
\setlength{\parindent}{0pt}
\setlength{\parskip}{6.5pt}
\setlength{\emergencystretch}{1.6em}
\setstretch{1.135}
\setlength{\JustifyingParindent}{0pt}


\begin{document}
% Coordonnees en quatre lignes comme dans la lettre de Matheo,
% en gardant la palette personnelle bleu et rouge.
{\setstretch{1.04}
{\fontsize{15.4}{17.2}\selectfont\bfseries\color{bleu}Youssef EL KATTOUFI}\par\vspace{1pt}
{\fontsize{10.1}{11.4}\selectfont\bfseries\color{rouge}Élève-ingénieur -- Électronique et informatique}\par\vspace{1pt}
{\fontsize{9.9}{11.1}\selectfont +33 6 51 65 42 37\quad\textcolor{gris}{|}\quad \href{mailto:elkattoufi.youssef@hotmail.com}{elkattoufi.youssef@hotmail.com}}\par
{\fontsize{9.9}{11.1}\selectfont Choisy-le-Roi, France}\par}
\vspace{20pt}

{\fontsize{9.35}{10.5}\selectfont\raggedleft
À l'attention de l'équipe Acquisition de talents Campus\\
Bell Canada\\[3pt]
Choisy-le-Roi, le 15 novembre 2026\par}
\vspace{14pt}
{\fontsize{10}{11.5}\selectfont\bfseries\color{bleu}Objet : Candidature au programme de stages 2027 -- Innovation technologique (réf. 432858)\par}
\vspace{4pt}

\fontsize{10.7}{12.35}\selectfont
Madame, Monsieur,

\justifying
Actuellement en quatrième année du cycle ingénieur en Électronique et Informatique à Polytech Sorbonne, je souhaite rejoindre Bell dans le cadre du programme de stages 2027 en innovation technologique. Je recherche un stage de quatre mois à partir d'avril 2027 et poursuivrai mes études à l'automne suivant.

J'ai eu l'occasion d'assister au webinaire de Bell consacré aux stages et aux nouveaux diplômés. J'ai notamment retenu que les stagiaires pouvaient contribuer à de vrais projets et que Bell développait aujourd'hui ses activités au-delà des télécommunications traditionnelles. C'est une évolution qui m'intéresse, car j'aimerais découvrir comment les différents métiers techniques travaillent ensemble sur des systèmes utilisés à grande échelle.

Je me suis également intéressé aux annonces concernant le Réseau d'IA tissé de Bell et son projet d'extension en Saskatchewan, ainsi qu'au déploiement, avec Cohere, d'une solution d'IA pour la cybersécurité. Ce qui attire mon attention, au-delà de l'IA, c'est tout ce qu'il faut faire fonctionner autour : les réseaux, les logiciels et les infrastructures. Cela rejoint mon intérêt pour l'électronique et l'informatique, et mon envie de participer à des projets qui ont une utilité concrète.

Je suis quelqu'un de calme. Lorsque je bloque, je commence par chercher des solutions de mon côté et par faire des essais. Je consulte ensuite des ressources sur Internet, puis je vais échanger avec les personnes autour de moi pour avoir d'autres points de vue. À une époque où l'IA peut proposer une réponse en une fraction de seconde, je trouve d'autant plus important de ne pas simplement la reprendre : j'essaie de comprendre la solution proposée, de la vérifier et de savoir pourquoi elle fonctionne, ou non.

J'aime aussi proposer mes idées, même si elles ne sont pas toujours les bonnes. J'aime en discuter avec les autres pour les améliorer. C'est ce que j'ai apprécié pendant la Coupe de France de robotique : j'ai développé des bases logicielles que les autres membres pouvaient compléter et je participais à la résolution des problèmes de programmation ou d'électronique. Cela m'a montré à quel point l'échange et le travail en équipe peuvent faire avancer un projet.

Chez Bell, je souhaiterais à la fois apprendre auprès d'ingénieurs plus expérimentés et apporter quelque chose à l'équipe : mes connaissances, ma persévérance et mes idées. Participer à des projets concrets, comprendre les choix techniques et contribuer à leur réalisation serait pour moi une expérience particulièrement enrichissante.

Je me tiens à votre disposition pour un entretien afin d'échanger davantage sur ma candidature et mes motivations.

Je vous prie d'agréer, Madame, Monsieur, l'expression de mes salutations distinguées.

\vspace{2pt}
{\raggedleft\bfseries\color{bleu}Youssef EL KATTOUFI\par}
\end{document}
